% GENERATED FILE. Edit canonical lessons, then run npm run generate:books.
% canonical-source-hash: fnv1a64:da186ae4da1d82ec
% canonical-lessons: PT-C378-rever, PT-C378-apontamentos, PT-C378-manual, PT-C378-decorar, PT-C378-teste

\chapter{Antes do exame}
\label{ch:pt-antes-do-exame}
\hlchaptermodality{378}

\begin{chapteropening}
\textbf{By the end of this chapter:} \emph{I can talk about getting ready for an exam: revising my notes and the textbook, learning by heart, and tests.}

Complete the last lesson of chapter 378: I can talk about getting ready for an exam: revising my notes and the textbook, learning by heart, and tests.
\end{chapteropening}

\section[rever]{rever — to revise, to go over; to see again}

\label{lesson:PT-C378-rever}



\begin{quote}
Before the new one: say the Portuguese for the master's degree, then the Portuguese for the doctorate.
\end{quote}

\subsection*{You'll want to know: rever}
\textbf{rever} — \textquotedblleft{}to revise, to go over; to see again\textquotedblright{}.

\begin{quote}
Revi a matéria toda no fim de semana.
\end{quote}

\begin{quote}
\emph{I revised all the material at the weekend.}
\end{quote}

Conjugated like \emph{ver}: \emph{revejo, revês, revê}. In Portugal it is the usual word for revising for exams.

\begin{cousinweb}
\emph{rever} is built from \emph{re-} \textquotedblleft{}again\textquotedblright{} and \emph{ver} \textquotedblleft{}to see\textquotedblright{}, inherited from Latin \emph{vidēre} \textquotedblleft{}to see\textquotedblright{}; compare Latin \emph{revidēre} \textquotedblleft{}to see again, to go to see again\textquotedblright{} (Lewis \& Short).
\end{cousinweb}

\begin{grammarlens}[title={What you've built}]
One more word for telling, explaining and describing at B1.
\end{grammarlens}

\subsection*{Your turn}
\begin{itemize}
\raggedright
  \item \emph{Say it:} \emph{rever}
  \item \emph{Say it:} \emph{rever}, once more
  \item \emph{Say it:} a sentence of your own with \emph{rever}
  \item \emph{Recall:} say the Portuguese for the master's degree, then the Portuguese for the doctorate, then say \emph{rever} again
\end{itemize}

\begin{tcolorbox}[breakable,colback=teal!4,colframe=teal!35!black,title={Before you move on}]
What does rever mean? (\textquotedblleft{}To revise, to go over; to see again\textquotedblright{}.) Say it once more.
\end{tcolorbox}
\section[os apontamentos]{os apontamentos — the class notes}

\label{lesson:PT-C378-apontamentos}



\begin{quote}
Before the new one: say the Portuguese for the doctorate, then the Portuguese for to revise.
\end{quote}

\subsection*{You'll want to know: os apontamentos}
\textbf{os apontamentos} — \textquotedblleft{}the class notes\textquotedblright{}.

\begin{quote}
Emprestei os meus apontamentos a uma colega.
\end{quote}

\begin{quote}
\emph{I lent my notes to a classmate.}
\end{quote}

European Portuguese; in Brazil \emph{as anotações}.

\begin{cousinweb}
\emph{apontamento} is built on \emph{apontar} \textquotedblleft{}to note down, to point\textquotedblright{}, on \emph{ponto} \textquotedblleft{}point\textquotedblright{}, inherited from Latin \emph{punctum} \textquotedblleft{}a point\textquotedblright{} (Lewis \& Short).
\end{cousinweb}

\begin{grammarlens}[title={What you've built}]
One more word for telling, explaining and describing at B1.
\end{grammarlens}

\subsection*{Your turn}
\begin{itemize}
\raggedright
  \item \emph{Say it:} \emph{os apontamentos}
  \item \emph{Say it:} \emph{os apontamentos}, once more
  \item \emph{Say it:} a sentence of your own with \emph{os apontamentos}
  \item \emph{Recall:} say the Portuguese for the doctorate, then the Portuguese for to revise, then say \emph{os apontamentos} again
\end{itemize}

\begin{tcolorbox}[breakable,colback=teal!4,colframe=teal!35!black,title={Before you move on}]
What does os apontamentos mean? (\textquotedblleft{}The class notes\textquotedblright{}.) Say it once more.
\end{tcolorbox}
\section[o manual]{o manual — the textbook; the manual}

\label{lesson:PT-C378-manual}



\begin{quote}
Before the new one: say the Portuguese for to revise, then the Portuguese for the class notes.
\end{quote}

\subsection*{You'll want to know: o manual}
\textbf{o manual} — \textquotedblleft{}the textbook; the manual\textquotedblright{}.

\begin{quote}
O manual de matemática pesa muito.
\end{quote}

\begin{quote}
\emph{The maths textbook is very heavy.}
\end{quote}

In Portugal, \emph{o manual escolar} is the school textbook. As an adjective, by hand: \emph{trabalho manual}.

\begin{cousinweb}
\emph{manual} is a learned borrowing from Latin \emph{manuālis} \textquotedblleft{}of the hand, that fills the hand\textquotedblright{} (EtymDB), built on \emph{manus} \textquotedblleft{}hand\textquotedblright{}: a book you hold.
\end{cousinweb}

\begin{grammarlens}[title={What you've built}]
One more word for telling, explaining and describing at B1.
\end{grammarlens}

\subsection*{Your turn}
\begin{itemize}
\raggedright
  \item \emph{Say it:} \emph{o manual}
  \item \emph{Say it:} \emph{o manual}, once more
  \item \emph{Say it:} a sentence of your own with \emph{o manual}
  \item \emph{Recall:} say the Portuguese for to revise, then the Portuguese for the class notes, then say \emph{o manual} again
\end{itemize}

\begin{tcolorbox}[breakable,colback=teal!4,colframe=teal!35!black,title={Before you move on}]
What does o manual mean? (\textquotedblleft{}The textbook; the manual\textquotedblright{}.) Say it once more.
\end{tcolorbox}
\section[decorar]{decorar — to learn by heart, to memorise; to decorate}

\label{lesson:PT-C378-decorar}



\begin{quote}
Before the new one: say the Portuguese for the class notes, then the Portuguese for the textbook.
\end{quote}

\subsection*{You'll want to know: decorar}
\textbf{decorar} — \textquotedblleft{}to learn by heart, to memorise; to decorate\textquotedblright{}.

\begin{quote}
Decorei o poema para a aula.
\end{quote}

\begin{quote}
\emph{I learnt the poem by heart for the class.}
\end{quote}

In Portugal \emph{decorar} very often means to memorise; it also means to decorate a room.

\begin{cousinweb}
\emph{decorar} \textquotedblleft{}to decorate\textquotedblright{} comes from Latin \emph{decorāre} \textquotedblleft{}to adorn\textquotedblright{} (EtymDB); the sense \textquotedblleft{}to learn by heart\textquotedblright{} is usually linked with \emph{de cor} \textquotedblleft{}by heart\textquotedblright{}, on Latin \emph{cor} \textquotedblleft{}heart\textquotedblright{} (Lewis \& Short).
\end{cousinweb}

\begin{grammarlens}[title={What you've built}]
One more word for telling, explaining and describing at B1.
\end{grammarlens}

\subsection*{Your turn}
\begin{itemize}
\raggedright
  \item \emph{Say it:} \emph{decorar}
  \item \emph{Say it:} \emph{decorar}, once more
  \item \emph{Say it:} a sentence of your own with \emph{decorar}
  \item \emph{Recall:} say the Portuguese for the class notes, then the Portuguese for the textbook, then say \emph{decorar} again
\end{itemize}

\begin{tcolorbox}[breakable,colback=teal!4,colframe=teal!35!black,title={Before you move on}]
What does decorar mean? (\textquotedblleft{}To learn by heart, to memorise; to decorate\textquotedblright{}.) Say it once more.
\end{tcolorbox}
\section[o teste]{o teste — the test}

\label{lesson:PT-C378-teste}



\begin{quote}
Before the new one: say the Portuguese for the textbook, then the Portuguese for to learn by heart.
\end{quote}

\subsection*{You'll want to know: o teste}
\textbf{o teste} — \textquotedblleft{}the test\textquotedblright{}.

\begin{quote}
Amanhã temos um teste de português.
\end{quote}

\begin{quote}
\emph{Tomorrow we have a Portuguese test.}
\end{quote}

At school, a written test during the term; \emph{o exame} is bigger.

\begin{cousinweb}
\emph{teste} was borrowed from English \emph{test} (EtymDB), which goes back through Old French \emph{test} \textquotedblleft{}pot for melting metal\textquotedblright{} to Latin \emph{testum} \textquotedblleft{}an earthen vessel\textquotedblright{} (Lewis \& Short): metals were tested in such pots.
\end{cousinweb}

\begin{grammarlens}[title={What you've built}]
One more word for telling, explaining and describing at B1.
\end{grammarlens}

\subsection*{Your turn}
\begin{itemize}
\raggedright
  \item \emph{Say it:} \emph{o teste}
  \item \emph{Say it:} \emph{o teste}, once more
  \item \emph{Say it:} a sentence of your own with \emph{o teste}
  \item \emph{Recall:} say the Portuguese for the textbook, then the Portuguese for to learn by heart, then say \emph{o teste} again
\end{itemize}

\begin{tcolorbox}[breakable,colback=teal!4,colframe=teal!35!black,title={Before you move on}]
What does o teste mean? (\textquotedblleft{}The test\textquotedblright{}.) Say it once more.
\end{tcolorbox}
